\begin{frame}{Credibility and trust}

\begin{alertblock}{\faShield~Industry threshold}
Credibility is fragile: early issues can undermine trust
\end{alertblock}

\begin{columns}[T,onlytextwidth]
\begin{column}{.50\textwidth}
\begin{block}{\faClipboardList~What is evaluated}
\begin{itemize}
    \item reproducibility
    \item transparency
    \item clear assumptions
\end{itemize}
\end{block}
\end{column}

\begin{column}{.45\textwidth}
\begin{exampleblock}{\faLightbulb~Shift in communication}
\centering
Build trust first \\
-- then invite inspection
\end{exampleblock}
\end{column}
\end{columns}

\end{frame}
\note{
As feedback moves from internal to external,  
the consequences change.  

In an earlier version of the course,  
students presented their work in a standard academic way.  

During one presentation, the industry partner identified  
an issue early in the analysis.  

Even though the rest of the work was reasonable,  
that early issue reduced trust in the entire result.  

So they did not engage further.  

This is what we mean by credibility being fragile.  

So we made a deliberate shift in how students communicate.  

The goal is to \emph{build trust first},  
so the audience is willing to engage more deeply.  

That means clearly presenting results, assumptions,  
and what can be done with the analysis.  

The detailed reasoning is still there --  
but supported through the report and the codebase.  

Because the full history is available,  
the work becomes \emph{traceable and reproducible}.  

So credibility is not just claimed --  
it can be inspected.
}